\SourceLeaf{264}%
\[
\frac{d\,dy}{dx^{2}} = -\alpha\alpha\,\frac{1-\dfrac{dy^{2}}{\alpha\alpha\,dx^{2}}}{y},
\]

\[
\frac{y\,ddy}{dx^{2}} - \frac{dy^{2}}{dx^{2}} + \alpha\alpha = 0.
\]

\noindent Mit $\dfrac{2\,dy}{y^{3}}$ multipliziert:

\[
\frac{2\dfrac{dy}{dx}\cdot d\dfrac{dy}{dx}}{yy} - \frac{2\left(\dfrac{dy}{dx}\right)^{2}dy}{y^{3}} + \frac{2\alpha\alpha\,dy}{y^{3}} = 0.
\]

\noindent Durch Integration

\[
\frac{1}{yy}\left(\frac{dy}{dx}\right)^{2} - \frac{\alpha\alpha}{yy} = \text{Const.} = -kk,
\]
\[
-kkyy + \alpha\alpha - \left(\frac{dy}{dx}\right)^{2} = 0,
\]
\[
dx = \frac{dy}{\BareRoot{(\alpha\alpha-kkyy)}} = \frac{\dfrac{k}{\alpha}dy}{k\BareRoot{\left(1-\dfrac{kk}{\alpha\alpha}yy\right)}},
\]
\[
kx = \text{Arc.}\sin\frac{k}{\alpha}y + \text{Const.},
\]
\[
y = \frac{\alpha}{k}\sin(kx+\lambda),
\]
\[
fx = \frac{\alpha}{k}\sin kx,
\]
\[
gx = \cos kx.
\]

\smallskip
\begin{center}

\end{center}

\begin{center}
\textsc{Bemerkungen.}
\end{center}

\noindent \textsc{Gauss} hat die vorstehende Notiz auf einem Zettel (4 Seiten 8\textsuperscript{0}) niedergeschrieben, der als Einlage in seinem Exemplar der \emph{»Geometrischen Untersuchungen zur Theorie der Parallellinien} von Nicolaus \textsc{Lobatschewsky}, Berlin 1840« gefunden wurde. Als Abfassungszeit darf daher 1840 bis 1846 angenommen werden.

Die Notiz soll unter der Voraussetzung, dass für Dreiecke mit drei unendlich kleinen Seiten die euklidische Geometrie gilt, die Gleichungen zwischen den Seiten und Winkeln eines endlichen Dreiecks herleiten. Es genügt, diese Beziehungen für ein rechtwinkliges Dreieck herzuleiten. Auf Grundlage des vorstehenden Textes lässt sich der Gedanke folgendermaßen ausführen.
